\subsection{Experiment 6B: acceleration and coordinate-consistency check}
\label{sec:experiment6b}

\paragraph{Question and scope.}
Experiment 6A supports the scheduling--specialization decomposition but leaves
an explicit prescribed-speed plant approximation. Does its principal conclusion
survive a consistent lateral-momentum model, and does it persist when the speed
profiles are slowed below the illustrative acceleration target? This is a
paired transfer check with the 6A fitted models and controller gains frozen.
It does not refit on corrected-plant data or introduce acceleration scheduling
into the controller.

\paragraph{Coordinate correction.}
For the small-angle sideslip proxy used throughout this benchmark,
$\beta=v_y/v_x$, the lateral-momentum and yaw equations give
\begin{align}
 \dot v_y&=\frac{F_{yf}+F_{yr}}{m}-v_xr,\\
 \dot r&=\frac{l_fF_{yf}-l_rF_{yr}}{I_z},\\
 \dot\beta&=\frac{F_{yf}+F_{yr}}{mv_x}-r
             -\frac{\dot v_x}{v_x}\beta.
 \label{eq:6b-coordinate}
\end{align}
The final term follows by differentiating $v_y/v_x$. It vanishes at constant
speed, but is absent from the previous prescribed-speed implementation.
We retain the same small-angle slip expressions and tanh force laws, written
in lateral-velocity coordinates as
\begin{align}
 \alpha_f&=\delta-\frac{v_y+l_fr}{v_x}, &
 \alpha_r&=\frac{-v_y+l_rr}{v_x}.
\end{align}
The controller still receives $(\beta,r)$ at each sample. This coordinate check
does not replace the proxy by the exact angle $\arctan(v_y/v_x)$.

\paragraph{Three plant implementations.}
The historical implementation, denoted \texttt{legacy}, retains the 6A
frozen-speed RK4 step. A diagnostic, \texttt{beta\_noacc}, varies speed between
RK4 stages while omitting the final term in \cref{eq:6b-coordinate}. The
coordinate-consistent implementation, \texttt{vy}, integrates $(v_y,r)$ and
converts $v_y$ to $\beta$ at sampling instants. Speed is linearly interpolated
between prescribed samples during integration. Steering remains held over each
sample, and the integral state and sampled controller are unchanged. Comparing
the first two implementations isolates within-step speed variation. Comparing
the latter two isolates the coordinate term, up to integration error.

All three retain static axle loads, pure lateral tanh tires, and externally
prescribed longitudinal speed. There is no longitudinal-force allocation,
combined-slip friction constraint, or acceleration-induced load transfer.
Thus \texttt{vy} is coordinate-consistent under the existing bicycle
assumptions, not a full longitudinal--lateral vehicle model.

\paragraph{Frozen protocol and contrasts.}
The four selected cells use $\gamma_h\in\{0,1\}$ and speed envelopes
18--22 and 10--30 m/s. We load the committed 6A parameter fits for seeds
21--30 and reconstruct exactly the same M1--M4 gains. Test fleets remain
the paired held-out realizations 121--130. These are repeated conditions on
the same clients, not additional independent fleet evidence. Per-seed fit-file
hashes record the model provenance.

Each of the three held-out scenarios and three severities is run with its
original 12-second timing and with both speed and yaw reference stretched to
24 seconds. The sample time remains 0.01 seconds. This retains reference
amplitudes and speed support while halving the rate of both signals. Broad-range
peak speed acceleration decreases from 3.93 to 1.96 m/s$^2$. Stretching the yaw
reference also makes lateral tracking less demanding, so we compare architectures
and plants within each timing condition rather than attributing cross-duration
differences solely to acceleration.

Primary metrics and amplitude-feasibility limits remain those of 6A. Paired
bootstrap intervals resample ten seeds, using 10,000 draws and seed 6201.
The predeclared robustness gate requires M4 to improve moderate mean tracking
by at least 5\% over both M2 and M3 in the nominal-separation broad-envelope
\texttt{vy} condition at both durations, with positive lower confidence bounds
and amplitude feasibility at all tested severities. The four-corner
smallest-controller classification is a secondary descriptive result using
the previous 5\% upper-bound rule. The 6A frozen-speed stability audit applies
to the unchanged gains and constant-speed linearizations, but does not certify
the accelerating trajectories.

\paragraph{Validation and reproducibility.}
Constant-speed closed-loop trajectories and metrics agree with the previous
simulator to numerical precision. The corrected derivative is checked against
the coordinate transformation, and RK4 is checked against an independently
implemented lateral-momentum ODE using adaptive integration and step halving.
The complete suite passes 33 tests. The full legacy replay is also compared
against the committed 6A raw metrics before generating conclusions.
The entry point is \texttt{experiment\_6b\_acceleration\_validation.py}, with
the frozen protocol under \texttt{code/config/}. Compressed per-seed client and
profile CSVs support paired summaries and a \texttt{--summarize-only} rebuild.

\paragraph{Results and decision.}
All 259,200 client evaluations complete and pass amplitude feasibility.
The maximum absolute discrepancy in the retained legacy-replay metrics is
$7.95\times10^{-13}$, confirming reproduction of the corresponding 6A runs.
The unchanged frozen audit has worst spectral radius 0.96888. The primary
robustness gate passes at both durations, with all four paired comparisons
favoring M4 in all ten fleet pairs and all lower confidence bounds positive.

\Cref{tab:6b-tracking} gives the nominal-separation broad-envelope results.
At the original timing, corrected-plant M4 improves on M3 by 16.46\% and on
M2 by 25.02\%. At doubled duration, the reductions are 35.27\% and 44.95\%,
respectively. The 95\% paired intervals for the absolute mean differences
$J_{M3}-J_{M4}$ and $J_{M2}-J_{M4}$ are $[0.00491,0.00580]$ and
$[0.00893,0.00922]$ degrees/s at 12 seconds, and $[0.00378,0.00447]$ and
$[0.00603,0.00634]$ degrees/s at 24 seconds.

\begin{table}[htbp]
\centering
\caption{Experiment 6B: mean moderate yaw-rate tracking RMSE in degrees/s for
nominal separation and 10--30 m/s. Models and gains are identical across plant
implementations and timings. The corrected plant integrates lateral momentum.}
\label{tab:6b-tracking}
\begin{tabular}{lrrrr}
\toprule
 & \multicolumn{2}{c}{12 seconds} & \multicolumn{2}{c}{24 seconds}\\
Method & Legacy & Corrected & Legacy & Corrected\\
\midrule
M1 & 0.040306 & 0.040336 & 0.016590 & 0.016597\\
M2 & 0.036234 & 0.036246 & 0.013748 & 0.013748\\
M3 & 0.032481 & 0.032533 & 0.011676 & 0.011692\\
M4 & 0.027164 & 0.027177 & 0.007561 & 0.007568\\
\bottomrule
\end{tabular}
\end{table}

The slower references reduce absolute error for all four methods. Although
M4's relative advantages increase, its absolute advantages over M2 and M3
decrease. This is not evidence that reducing acceleration itself increases the
benefit of LPV. Both the lateral reference rate and its interaction with speed
change in this timing experiment.

\Cref{fig:6b-acceleration} shows that the three plant implementations nearly
overlap at the scale of the architecture differences. Across the four cells,
two durations, and four methods, the largest absolute relative change in
moderate mean tracking from \texttt{legacy} to \texttt{vy} is 0.1594\%.
Changing within-step speed treatment alone produces at most 0.01394\% change.
For broad-range nominal-family M4, the complete plant correction increases
mean error by 0.0484\% at 12 seconds and 0.0938\% at 24 seconds. These changes
are small but statistically resolvable in the paired simulations. They are
not zero, and are much smaller than the architecture contrasts in this study.

\begin{figure}[htbp]
\centering
\includegraphics[width=0.95\linewidth]{experiment_6b_acceleration.pdf}
\caption{Experiment 6B: nominal-separation broad-range tracking across plant
implementations. Error bars are standard deviations across ten seed pairs.
\texttt{beta\_noacc} varies speed within RK4 but omits the coordinate term,
whereas \texttt{vy} integrates the consistent momentum equations. Their near
overlap shows that the controller-architecture differences dominate these
plant corrections under both tested timings.}
\label{fig:6b-acceleration}
\end{figure}

The four-corner descriptive classifications also remain unchanged at both
durations and for all three implementations. Global LTI is sufficient with
zero separation and narrow range, global LPV with zero separation and broad
range, family LTI with nominal separation and narrow range, and family LPV
with nominal separation and broad range. These are the corresponding four
corners of \cref{tab:6a-regimes}. The intermediate-separation and
intermediate-envelope boundaries of 6A were not rerun here.

One reason a modest effect is plausible is the size of the added term relative
to linear-tire sideslip damping. Their coefficient ratio is
\begin{equation}
 \frac{|\dot v_x|/v_x}{(C_f+C_r)/(mv_x)}
 =\frac{m|\dot v_x|}{C_f+C_r}.
\end{equation}
At the nominal center, $(C_f+C_r)/m\simeq106.7\,\mathrm{m\,s^{-2}}$, so
even the original peak speed acceleration gives approximately 3.7\% of that
linear damping coefficient. This is an explanatory scale comparison, not an
error bound for tanh tires or the closed loop. The small observed tracking
effect is established by the paired simulations.

The result supports continuing toward limited-data identification and
federation: the tested M4 advantage is not an artifact of the omitted
coordinate term or dependent on the faster timing. The remaining limitations
are the prescribed longitudinal speed, pure lateral tire model, fixed
controllers, and absence of a time-varying stability certificate. The result
does not yet validate identification on corrected-plant data or the 5A
controller-grid sweep under the new plant.

\paragraph{Implication for subsequent identification.}
The coordinate term is known kinematics shared by all clients, not a new
structural family. Identification on accelerating trajectories should either
use lateral-velocity coordinates or include the known term
$-(\dot v_x/v_x)\beta$ in its predictor. The persistent unknown parameter ratios
can then retain their original interpretation. This experiment tests transfer
of existing speed-only controllers and does not claim that the previous
identification objective is already corrected. Estimating acceleration from
noisy production signals remains a separate practical question.
